\chapter{系统总体设计}

\section{系统需求分析}

本节填写课程设计的需求分析内容。可根据任务书要求，从功能需求、性能需求、硬件资源需求、软件设计需求和测试验证需求等方面进行说明。需求分析应回答“系统需要实现什么功能”“需要满足什么条件”“最终如何判断设计完成”等问题。

功能需求部分可说明系统需要完成的主要功能，例如输入采集、数据处理、状态控制、结果显示、通信调试或报警提示等。不同课题可根据实际内容进行取舍，不必机械套用。

性能需求部分可说明系统在稳定性、实时性、准确性、可靠性、资源占用或可扩展性等方面的要求。若课程设计没有明确量化指标，也可以从工程实现角度进行定性说明。

测试验证需求部分可说明后续需要通过哪些实验现象或测试数据验证系统功能，例如功能是否正常、状态切换是否正确、显示结果是否符合预期、异常情况是否能够处理等。

\section{系统总体方案}

本节填写系统总体方案设计内容。可从系统组成、模块划分、数据流向和工作流程等方面进行说明，明确系统由哪些部分组成，各部分之间如何连接或协同工作。

系统总体方案一般包括输入模块、核心处理模块、输出显示模块、通信调试模块、电源模块或其他扩展模块。实际撰写时应结合具体课题内容进行调整，不需要包含与本设计无关的模块。

系统总体方案如图~\ref{fig:system_overview} 所示。该图应展示系统的主要组成部分及其相互关系，为后续硬件设计和软件设计提供总体框架。

\begin{figure}[htbp]
    \centering
    \includegraphics[width=0.8\textwidth]{fig_2_1_system_overview}
    \caption{系统总体方案框图}
    \label{fig:system_overview}
\end{figure}

\section{硬件平台与主控选型}

本节填写硬件平台与主控芯片选型内容。可介绍课程设计所使用的开发板、主控芯片、主要外设资源以及选择该平台的原因。

主控芯片选型应结合课题功能需求进行说明，例如处理能力、片上外设资源、接口数量、开发生态、程序下载与调试便利性等。重点说明所选平台为什么能够满足本课程设计的功能实现和测试验证要求。

若系统包含传感器、显示模块、通信模块、执行模块或其他外设，也可在本节简要说明其选型依据。更详细的引脚连接和硬件资源分配可放在第三章展开。

\section{关键技术方案}

本节填写本课程设计涉及的关键技术方案。可根据具体课题内容选择说明一种或多种关键技术，例如信号采集方案、数据处理方案、控制逻辑方案、显示交互方案、通信调试方案或状态机设计方案等。

关键技术方案的描述应突出“为什么这样设计”和“该方案如何支撑系统功能”。不需要在本节展开所有程序细节，具体的软件流程、核心函数和实现逻辑可放在第四章说明。

如果本课题涉及特定检测原理、控制算法、通信协议或显示刷新机制，可在本节进行原理性说明，并为后续章节的硬件连接和软件实现做好铺垫。

\section{系统工作流程}

本节填写系统整体工作流程。可按照系统上电初始化、外设配置、主循环运行、输入检测、逻辑判断、输出更新和测试反馈等步骤进行描述。

系统工作流程一般可概括为：首先完成硬件外设和软件变量初始化；随后进入主循环或任务调度流程；在运行过程中周期性采集输入信息，完成数据处理和状态判断；最后根据判断结果更新输出模块或显示模块。


\section{本章小结}

本章主要对课程设计系统进行了总体设计说明，包括需求分析、总体方案、硬件平台与主控选型、关键技术方案和系统工作流程等内容。通过本章设计，明确了系统的整体组成和实现思路，为后续硬件设计、软件设计及系统测试奠定了基础。
